\documentclass[12pt]{article}
\usepackage[legalpaper,
            lmargin=1in,rmargin=0.5in,
            bmargin=1in,tmargin=1in]{geometry}
\usepackage{fontspec}
\usepackage[utf8]{inputenc}
\usepackage{titlesec}
\usepackage{multirow}
\usepackage{graphicx}
\usepackage{titling}
\usepackage{array}
\usepackage{makecell}
\usepackage{fancyhdr}
\usepackage[none]{hyphenat}
\usepackage{longtable}
\usepackage{adjustbox}
\usepackage[dvipsnames]{xcolor}
\usepackage[]{amsmath}
\usepackage[banglamainfont=Kalpurush,
            banglattfont=SolaimanLipi,
            % feature=1,
            % changecounternumbering=0
           ]{latexbangla}

\newcommand{\my}{\includegraphics[height=3.2em]{pic/my.png}}
\pagestyle{fancy}
\fancyhf{}
\renewcommand{\headrulewidth}{0pt}
% header
\lhead{
\framebox[1.1\width]{\footnotesize{``উন্নয়নের অক্সিজেন রাজস্ব''}}
}
\rhead{
\my
\\
\framebox[1.1\width]{\footnotesize{``জনকল্যানে রাজস্ব''}}
}
% footer

\fancypagestyle{laststyle}
{
   \fancyhead[C]{\underline{পাতা নং-}}
   \fancyhead[L]{}
   \fancyhead[R]{}
}

% \newcommand{\fileno}{নথি নং- এস-৫-১৬৩/স্টাফ/কাস/০৬-০৭/অংশ-২}
\newcommand{\fileno}{নথি নং- \hspace{5em}/স্টাফ/কাস/০৬-০৭/অংশ-  }
\newcommand{\filenou}{
\begin{minipage}[t]{0.60\linewidth}
\hspace{1em}
\end{minipage}
\begin{minipage}[t]{0.40\linewidth}
\underline{\fileno}
\end{minipage}
}
\newcommand{\nfpage}{
\newpage
\small
{\filenou}
}
\newcommand{\fdt}{\hspace*{3em} তারিখ: \hspace{2.4em} /১১/২১}
\newcommand{\frbdt}{তারিখ: \hspace{2.4em} /১১/২১}
\newcommand{\pack}{STC:1245 BALES}
\newcommand{\good}{UNITED STATES RAW COTTON}
\newcommand{\co}{UNITED STATES}
\newcommand{\depo}{SAPL, DEPOT, CHITTAGONG}
\newcommand{\rdepo}{..................................., DEPOT, CHITTAGONG}
\newcommand{\vessel}{HR AARAI}
\newcommand{\rotno}{2021/5040}
\newcommand{\blno}{MEDUU2181458}
\newcommand{\beno}{C-1743011}
\newcommand{\bedt}{31.10.2021}
\newcommand{\rno}{\hspace{7em} / \hspace{4em}}
\newcommand{\rdt}{\hspace{2em}-11-21}
\newcommand{\conno}{$\boldsymbol{14}$}
\newcommand{\conn}{] FCL Container}
\newcommand{\cond}{$\boldsymbol{\times 40^\prime}$}
% \newcommand{\cond}{$\boldsymbol{\times 40^\prime\hspace{1em}\&\hspace{1em}05\times 40^\prime}$}
\newcommand{\ssml}{SHAMEEM SPINNING MILLS LTD.}
\newcommand{\scml}{SHAMEEM COMPOSITE MILLS LTD.}
\newcommand{\jsml}{JAMUNA SPINNING MILLS LTD.}
\newcommand{\jsmla}{SHAFIPUR, KALIAKAIR, GAZIPUR-1751, BANGLADESH.}
\newcommand{\jeal}{JAMUNA ELECTRONICS \& AUTOMOBILES LTD.}
\newcommand{\jdl}{JAMUNA DENIMS LTD.}
\newcommand{\jdla}{SHAFIPUR, KALIAKAIR, GAZIPUR-1751, BANGLADESH.}
\newcommand{\jdwl}{JAMUNA DENIMS \& WEAVING LTD.}
\newcommand{\htf}{HOORAIN HTF LTD.}
\newcommand{\impn}{\jdl}
\newcommand{\impadd}{\jdla}
\newcommand{\cnfn}{SHAMEEM SPINNING MILLS LTD.}
\newcommand{\cnfadd}{92,HIGH LEVEL ROAD, LALKHAN BAZAR, CHITTAGONG}
\newcommand{\ssmlreg}{২৩১৩৯৫০৬১৮০-এইচ, তারিখ: ০৯/০১/২০১৩}
\newcommand{\jealreg}{এল-২৯৩০১৩০১৩৪৩১-এইচ, তারিখ: ০৯/০১/২০১৩}
\newcommand{\impreg}{\hspace{10em} তারিখ: \hspace{5em}}
\newcommand{\continental}{CONTINENTAL TRADERS (BD) LIMITED}
\newcommand{\continentala}{IIUC Tower (9th floor), Plot \#09, Agrabad C/A, Sk. Mujib Road, Chittagong-4100.}
\newcommand{\transmarine}{TRANSMARINE LOGISTICS LIMITED.}
\newcommand{\transmarinea}{B.M Heights(4th Flr),318,SK Mujib Road,Agrabad C/A, Chattogram.}
\newcommand{\mscmed}{MSC MEDITERRANEAN SHIP.CO.BD.LTD.}
\newcommand{\mscmeda}{IIUC TOWER 4TH FLOOR, 1700/A SK.MUJIB ROAD PLOT-09, AGRABAD, CHITTAGONG}
\newcommand{\san}{\mscmed}
\newcommand{\sad}{\mscmeda}
\newcommand{\rodt}{\hspace*{3.0em} তারিখ: \hspace{3.5em}/১১/২০২১ ইং}
\newcommand{\robdt}{\hspace*{3.0em} \hspace{3.5em}/১১/২০২১ ইং}
\newcommand{\consh}{
\noindent{
\begin{tabular}{|l|l|l|l|l|}
\hline
\textbf{BMOU3084795}{\cond}
&
\textbf{TTNU5702348} {\cond}
&
\textbf{MSCU4760409} {\cond}
&
\textbf{MSDU7726021} {\cond}
&
\textbf{MSDU4004977} {\cond}
\\
\hline
\textbf{MEDU4132245}{\cond}
&
\textbf{MSDU4062063} {\cond}
&
\textbf{MEDU4002771} {\cond}
&
\textbf{MEDU4295860} {\cond}
&
\textbf{MEDU4172952} {\cond}
\\
\hline
\textbf{MSCU5969353}{\cond}
&
\textbf{DRYU4183999} {\cond}
&
\textbf{MSCU4490915} {\cond}
&
\textbf{FCIU8712660} {\cond}
&
\\
\hline
\end{tabular}
}
}

\begin{document}
\begin{center}
\vspace*{0.2\baselineskip}
\footnotesize{গণপ্রজাতন্ত্রী বাংলাদেশ সরকার}
\\
\footnotesize{কাস্টম হাউস, চট্টগ্রাম।}
\end{center}
\begin{minipage}[t]{.76\linewidth}
\hspace{0.5em}
\end{minipage}
\begin{minipage}[t]{0.24\linewidth}
\footnotesize{রিস্ক বন্ড নং - .....................}
\end{minipage}
\begin{minipage}[t]{.74\linewidth}
\footnotesize{{\fileno}}
\end{minipage}
\begin{minipage}[t]{.26\linewidth}
\footnotesize{{\fdt}}
\\
\end{minipage}
\footnotesize{বরাবরে
\\
SHAMEEM SPINNING MILLS LTD.
\\
92, HIGH LEVEL ROAD,
\\
LALKHAN BAZAR,
\\
CHITTAGONG.}
\\
\\
\begin{minipage}[t]{0.04\linewidth}
\footnotesize{বিষয়:}
\end{minipage}
\begin{minipage}[t]{0.96\linewidth}
\textbf{\footnotesize{আমদানিকৃত
[{\pack}, {\good}]
C/O. {\co}
হতে মালামালসহ
% container
[{\conno}{\cond}{\conn}
চট্টগ্রাম বন্দরের আওতাধীন
{\depo} হতে
আমদানিকারকের কারখানা চত্বরে
নিয়ে যাওয়ার অনুমতি প্রসঙ্গে।}}
\\
\footnotesize{Vessel: {\vessel},
ROT NO. {\rotno},
B/L NO. \textbf{{\blno}}
\\
B/E NO. \textbf{\beno},
DT: {\bedt},
BANK/TREASURY NO: {\rno},
DT: {\rdt}}
\\
\footnotesize{
[{\conno}{\cond}{\conn}}
\\
\scriptsize{
{\consh}
}
\\
\footnotesize{
\framebox[1.02\width]{IMP:{\impn}}
\\
\framebox[1.02\width]{FACTORY:{\impadd}}}
\\
\end{minipage}
\footnotesize{উপর্যুক্ত বিষয়ে আপনার পত্র নং-নাই,
তারিখ: \hspace{3em}/১০/২০২১ ইং
এর প্রেক্ষিতে জানানো যাচ্ছে যে, সংশ্লিষ্ট
\textbf{Shipping Agent: {\san}},
{\sad}
কর্তৃক অনাপত্তি প্রদান করায় এবং সংশ্লিষ্ট
শুল্কায়ন সেকশন কর্তৃক পণ্যচালানটি শুল্কায়নের
বিপরীতে ইস্যুকৃত রিলিজ অর্ডারের ভিত্তিতে উপরোক্ত
% container
[{\conno}{\cond}{\conn}
এ আমদানিকৃত মালামালসহ আমদানিকারকের নিজস্ব
কারখানা চত্বর {\impadd} -এ নিম্নোক্ত
শর্তে খালাস করার অনুমতি
প্রদান করা হলো।}
\\
\begin{minipage}[t]{0.04\linewidth}
\hspace{0.5em}
\end{minipage}
\begin{minipage}[t]{0.04\linewidth}
\scriptsize{(ক)}
\end{minipage}
\begin{minipage}[t]{.92\linewidth}
\scriptsize{আমদানিকৃত আলোচ্য চালানের পণ্য
কন্টেইনারসহ বেসরকারি ডিপো
{\depo}
হতে স্থানান্তর কালে স্থায়ী
আদেশ নং-৯৬, তারিখ: ১০/১২/২০১৪ ইং
মোতাবেক অনুসরণে অফ ডকে
কর্মরত কাস্টম কর্মকর্তা, রাজস্ব কর্মকর্তা,
ডিপোতে এআরও/আরও কর্তৃক পদস্থ সহকারী রাজস্ব
কর্মকর্তা এবং ডিপোতে নিয়োজিত সহকারী
রাজস্ব কর্মকর্তা উল্লিখিত Container
নম্বর সঠিক প্রাপ্তি সাপেক্ষে কন্টেইনার খুলে
ঘোষিত মালামাল পরীক্ষাপূর্বক সঠিক পেলে শুল্কায়িত
বি/ই ও রিলিজ অর্ডারের আলোকে ডেলিভারী অনুমতি
দিবেন। কোন প্রকার গড়মিল পরিলক্ষিত হলে
কন্টেইনারসহ মালামাল ডেলিভারী কার্যক্রম স্থগিত রেখে
তাৎক্ষনিকভাবে তা উর্দ্ধতন কর্তৃপক্ষের গোচরীভূত করবেন।}
\\
\end{minipage}
\begin{minipage}[t]{0.04\linewidth}
\hspace{0.5em}
\end{minipage}
\begin{minipage}[t]{0.04\linewidth}
\scriptsize{(খ)}
\end{minipage}
\begin{minipage}[t]{0.92\linewidth}
\scriptsize{আমদানিকারকের কারখানা চত্বরে
মালামাল খালাসের পর খালি কন্টেইনার পণ্য খালাসের
৩০ (ত্রিশ) দিনের মধ্যে সংশ্লিষ্ট শিপিং এজেন্ট কর্তৃক
প্রদত্ত অনাপত্তির শর্ত মোতাবেক স্থায়ী
আদেশ নং-১৩, তারিখ: ০৪/১১/০৩ ইং
অনুযায়ী শিপিং এজেন্টের মনোনীত কন্টেইনার
ইয়ার্ড/ডিপো {\rdepo} এ ফেরৎ এনে সংশ্লিষ্ট
শিপিং এজেন্ট, কাস্টম কর্তৃপক্ষের মাধ্যমে হস্তান্তর করে
স্থানান্তরের প্রমাণস্বরূপ ডুপ্লিকেট রিস্ক বন্ডের পিছনে
সংশ্লিষ্ট সকল পক্ষের গৃহীত স্বাক্ষরসহ
(রিস্ক বন্ড অবমুক্ত করার জন্য)
স্টাফ শাখায় দাখিল করবেন।
তাছাড়া রাজস্ব কর্মকর্তা সংশ্লিষ্ট (ডিপো) কর্তৃক
স্বাক্ষরিত প্রতিবেদনের কপি সংযুক্ত করতে হবে।
ইহার ভিত্তিতে আমদানিকারক কর্তৃক দেয়
রিস্ক বন্ড বাতিল করা হবে।}
\\
\\
\\
\\
\\
\\
\end{minipage}
\begin{minipage}[t]{0.72\linewidth}
\hspace{1em}
\end{minipage}
\begin{minipage}[t]{0.28\linewidth}
\hspace*{3.4em}\footnotesize{(এ.এম.নাহিদ হোসেন)}
\\
\hspace*{3.4em}\footnotesize{রাজস্ব কর্মকর্তা (স্টাফ)}
\\
\scriptsize{সহকারী/ডেপুটি কমিশনার অব কাস্টমস এর পক্ষে}
\\
\hspace*{5.5em}\scriptsize{কাস্টম হাউস, চট্টগ্রাম।}
\\
\scriptsize{{\rodt}}
\end{minipage}
\\
\scriptsize{{\fileno}}
\\
\underline{\scriptsize{অনুলিপি অবগতি ও প্রয়োজনীয় ব্যবস্থা গ্রহনের জন্য প্রেরণ করা হলো}}
\\
\begin{minipage}[t]{0.06\linewidth}
\scriptsize{০১।}
\end{minipage}
\begin{minipage}[t]{0.94\linewidth}
\scriptsize{বিভাগীয় কর্মকর্তা, কাস্টমস এক্সাইজ ও ভ্যাট, (গাজীপুর বিভাগ),
বাড়ি নং- ৩১৯,
খন্দকার কমপ্লেক্স, নলজানী,
চান্দানা চৌরাস্তা, জয়দেপপুর রোড,
গাজীপুর,
পন্যচালানটি উপরোক্ত প্রতিষ্ঠানে যথাযথভাবে
ব্যবহার করা হচ্ছে/হয়েছে কিনা সে
সম্পর্কিত একটি প্রতিবেদন নিম্ন স্বাক্ষরকারীর
দপ্তরে প্রেরণ করার জন্য সংশ্লিষ্ট সার্কেলকে
প্রয়োজনীয় নির্দেশ দানের জন্য আপনাকে
বিশেষভাবে অনুরোধ করা হলো।}
\end{minipage}
\begin{minipage}[t]{0.06\linewidth}
\scriptsize{০২।}
\end{minipage}
\begin{minipage}[t]{0.94\linewidth}
\scriptsize{ইপপোর্ট ম্যানেজার {\depo}}
\end{minipage}
\begin{minipage}[t]{0.06\linewidth}
\scriptsize{০৩।}
\end{minipage}
\begin{minipage}[t]{0.94\linewidth}
\scriptsize{রাজস্ব কর্মকর্তা {\depo}}
\end{minipage}
\begin{minipage}[t]{0.06\linewidth}
\scriptsize{০৪।}
\end{minipage}
\begin{minipage}[t]{0.94\linewidth}
\scriptsize{রাজস্ব কর্মকর্তা, সংশ্লিষ্ট গেইট ডিভিশন, {\depo}}
\end{minipage}
\begin{minipage}[t]{0.06\linewidth}
\scriptsize{০৫।}
\end{minipage}
\begin{minipage}[t]{0.94\linewidth}
\scriptsize{সহকারী রাজস্ব কর্মকর্তা, {\rdepo}}
\end{minipage}
\begin{minipage}[t]{0.06\linewidth}
\scriptsize{০৬।}
\end{minipage}
\begin{minipage}[t]{0.94\linewidth}
\scriptsize{শিপিং এজেন্ট:
{\san}, {\sad}}
\end{minipage}
\begin{minipage}[t]{0.06\linewidth}
\scriptsize{০৭।}
\end{minipage}
\begin{minipage}[t]{0.94\linewidth}
\scriptsize{IMP: {\impn}
FACTORY:{\impadd}}
\end{minipage}
\begin{minipage}[t]{0.06\linewidth}
\scriptsize{০৮।}
\end{minipage}
\begin{minipage}[t]{0.94\linewidth}
\scriptsize{অফিস কপি।}
\\
\\
\\
\\
\\
\\
\end{minipage}
\begin{minipage}[t]{0.72\linewidth}
\hspace{1em}
\end{minipage}
\begin{minipage}[t]{0.28\linewidth}
\scriptsize{{\robdt}}
\\
\hspace*{4.2em}\footnotesize{(এ.এম.নাহিদ হোসেন)}
\\
\hspace*{3.4em}\footnotesize{রাজস্ব কর্মকর্তা (স্টাফ)}
\\
\scriptsize{সহকারী/ডেপুটি কমিশনার অব কাস্টমস এর পক্ষে}
\\
\hspace*{5.5em}\scriptsize{কাস্টম হাউস, চট্টগ্রাম।}
\end{minipage}
{\nfpage}
\normalsize
\begin{minipage}[t]{0.04\linewidth}
।
\end{minipage}
\begin{minipage}[t]{0.96\linewidth}
বিষয়:\hspace{1em}\textbf{আমদানিকৃত
[{\pack}, {\good}]
C/O. {\co}
হতে মালামালসহ
% container
[{\conno}{\cond}{\conn}
চট্টগ্রাম বন্দরের আওতাধীন
{\depo} হতে
আমদানিকারকের কারখানা চত্বরে
নিয়ে যাওয়ার অনুমতি প্রসঙ্গে।}
\\
\footnotesize{Vessel: {\vessel},
ROT NO. {\rotno},
B/L NO. \textbf{{\blno}}
\\
B/E NO. \textbf{\beno},
DT: {\bedt},
BANK/TREASURY NO: {\rno},
DT: {\rdt}}
\\
% container
\footnotesize{
[{\conno}{\cond}{\conn}}
\\
\scriptsize{
{\consh}
}
\\
\footnotesize{
\framebox[1.02\width]{IMP:{\impn}}
\\
\framebox[1.02\width]{FACTORY:{\impadd}}
\\
\framebox[1.02\width]{ছাড়কারী:{\cnfn},{\cnfadd}}}
\\
\end{minipage}
\begin{minipage}[t]{0.04\linewidth}
\hspace{0.0em}
\end{minipage}
\begin{minipage}[t]{0.96\linewidth}
\normalsize
উপরোক্ত পত্র সংযোজনীসহ নথির পত্রাংশে উপস্থাপিত।
অনুগ্রহ পূর্বক দেখা যেতে পারে।
উপরোক্ত আমদানিকারক কর্তৃক আমদানিকৃত
[{\conno}{\cond}{\conn}
এর মাধ্যমে আনীত মালামাল সহ
আমদানিকারকের কারখানা চত্বর
{\impadd}-এ
নিয়ে খালাস করার অনুমতি চেয়ে
সিএন্ডএফ এজেন্ট: {\cnfn}, {\cnfadd}
আবেদন করেছেন।
\\
\end{minipage}
\footnotesize
\begin{minipage}[t]{0.04\linewidth}
।
\end{minipage}
\begin{minipage}[t]{0.96\linewidth}
আমদানিকৃত মালামাল সহ কন্টেইনার ডেলিভারী
সংক্রান্ত বিদ্যমান স্থায়ী আদেশ নং-০৮,
তারিখ: ২৯/০৭/০৩ খ্রি: এর আলোকে
আমদানিকারকের অত্র প্রতিষ্ঠানটি বিনিয়োগ
বোর্ডের আওতায় নিবন্ধিত একটি শিল্প
প্রতিষ্ঠান। যার নিবন্ধন সংখ্যা {\impreg}।
উক্ত স্থায়ী আদেশের আলোকে
নিম্নোক্ত শর্তে কন্টেইনার ডেলিভারীর
অনুমতি প্রদান করা যেতে পারে।
\\
\begin{minipage}[t]{0.04\linewidth}
(ক)
\end{minipage}
আমদানিকৃত আলোচ্য চালানের পণ্য
কন্টেইনারসহ বেসরকারি ডিপো
{\depo}
হতে স্থানান্তর কালে স্থায়ী
আদেশ নং-৯৬, তারিখ: ১০-১২-২০১৪ ইং
মোতাবেক অনুসরণে অফ ডকে
কর্মরত কাস্টম কর্মকর্তা, রাজস্ব কর্মকর্তা,
ডিপোতে এআরও/আরও কর্তৃক পদস্থ সহকারী রাজস্ব
কর্মকর্তা এবং ডিপোতে নিয়োজিত সহকারী
রাজস্ব কর্মকর্তা উল্লিখিত Container
নম্বর সঠিক প্রাপ্তি সাপেক্ষে কন্টেইনার খুলে
ঘোষিত মালামাল পরীক্ষাপূর্বক সঠিক পেলে শুল্কায়িত
বি/ই ও রিলিজ অর্ডারের আলোকে ডেলিভারী অনুমতি
দিবেন। কোন প্রকার গড়মিল পরিলক্ষিত হলে
কন্টেইনারসহ মালামাল ডেলিভারী কার্যক্রম স্থগিত রেখে
তাৎক্ষনিকভাবে তা উর্দ্ধতন কর্তৃপক্ষের গোচরীভূত করবেন।
\\
\end{minipage}
\begin{minipage}[t]{0.04\linewidth}
।
\end{minipage}
\begin{minipage}[t]{0.96\linewidth}
আমদানিকারকের কারখানা চত্বরে মালামাল
খালাসের পর খালি কন্টেইনার পণ্য খালাস
পরবর্তী ৩০ (ত্রিশ) দিনের মধ্যে
সংশ্লিষ্ট শিপিং এজেন্টের নিজস্ব কন্টেইনার
ইয়ার্ড/ডিপোট {\rdepo} -এ ফেরত প্রদান করবেন
মর্মে আমাদানিকারক কর্তৃক রিস্ক বন্ড দাখিল
করা হয়েছে।
সংশ্লিষ্ট মালামালের শুল্কায়িত বি/ই,
রিলিজ অর্ডার, এ্যাসেসমেন্ট নোটিশ,
ইনভয়েস, বি/এল, প্যাকিং লিস্ট,
ভ্যাট রেজিস্ট্রেশন, বিনিয়োগ বোর্ডের
পত্রের ফটোকপি ইত্যাদি দাখিল করা হয়েছে।
\\
\end{minipage}
\begin{minipage}[t]{0.04\linewidth}
।
\end{minipage}
\begin{minipage}[t]{0.96\linewidth}
সিএন্ডএফ এজেন্টের আবেদনের প্রক্ষিতে
সংশ্লিষ্ট শিপি এজেন্ট:
{\san}, {\sad}
কর্তৃক কন্টেইনার স্থানন্তরের কোন আপত্তি নেই
মর্মে অনাপত্তি প্রদান করা হয়েছে।
\\
\end{minipage}
\begin{minipage}[t]{0.04\linewidth}
।
\end{minipage}
\begin{minipage}[t]{0.96\linewidth}
আমদানিকারকের কারখানা চত্বরে মালামাল
খালাসের পর খালি কন্টেইনার পণ্য খালাস পরবর্তী
৩০ (ত্রিশ) দিনের মধ্যে
সংশ্লিষ্ট শিপিং এজেন্ট কর্তৃক প্রদত্ত
অনাপত্তির শর্ত মোতাবেক
স্থায়ী আদেশ নং-১৩,
তারিখ: ০৪/১১/০৩ ইং
অনুযায়ী শিপিং এজেন্টের মনোনীত
কন্টেইনার ডিপো/ইয়ার্ড {\rdepo}-এ
ফেরৎ এনে সংশ্লিষ্ট শিপিং এজেন্ট,
কাস্টম কর্তৃপক্ষের মাধ্যমে হস্তান্তর করে
স্থানান্তরের প্রমাণস্বরূপ ডুপ্লিকেট রিস্ক বন্ডের পিছনে
সংশ্লিষ্ট সকল পক্ষের গৃহীত স্বাক্ষরসহ
(রিস্ক বন্ড অবমুক্ত করার জন্য)
স্টাফ শাখায় দাখিল করবেন।
তাছাড়া আদেশ মোতাবেক
সহকারী/রাজস্ব কর্মকর্তা সংশ্লিষ্ট (ডিপো) কর্তৃক
স্বাক্ষরিত প্রতিবেদনের কপি সংযুক্ত করতে হবে।
ইহার ভিত্তিতে আমদানিকারক কর্তৃক দেয়
রিস্ক বন্ড বাতিল করা হবে।
\\
\end{minipage}
\begin{minipage}[t]{0.04\linewidth}
।
\end{minipage}
\begin{minipage}[t]{0.96\linewidth}
উল্লেখ্য অত্র আমদানিকারকের অনুকূলে ইতিপূর্বে
অনুমোদিত রিস্ক বন্ডের মেয়াদ উত্তীর্ণ কোন
রিস্ক বন্ড বাতিলের অপেক্ষায় অবশিষ্ট নাই।
\\
\end{minipage}
\begin{minipage}[t]{0.04\linewidth}
।
\end{minipage}
\begin{minipage}[t]{0.96\linewidth}
বিনিয়োগ বোর্ডের প্রত্যয়ন পত্র,
ভ্যাট রেজিষ্ট্রেশন কপি,
আউটপাশ বি/ই এবং প্রেরিতব্য পত্রে
আমদানিকারকের প্রতিষ্ঠানের কারখানা চত্বরের
ঠিকানা সঠিক আছে।
\\
\end{minipage}
\begin{minipage}[t]{0.04\linewidth}
।
\end{minipage}
\begin{minipage}[t]{0.96\linewidth}
এমতাবস্থায় আমদানীকারকের পক্ষে
পেশকৃত আবেদনপত্রের অনুকূলে
উপস্থাপিত রিস্ক বন্ডটি গ্রহণ করা যেতে পারে।
উল্লেখ্য, রিস্ক বন্ডে-২,০০,০০০/- টাকার
বিপরীতে ৫০০/- টাকার নন-জুডিশিয়াল
স্ট্যাম্পে রিস্ক বন্ড নং ..........................
{\frbdt} দাখিল করেছেন।
\\
\end{minipage}
\begin{minipage}[t]{0.04\linewidth}
।
\end{minipage}
\begin{minipage}[t]{0.96\linewidth}
বিনিয়োগ বোর্ডের নিবন্ধন পত্র এবং
মূসক ২.৩ নথিতে সংযুক্ত আছে।
অবমুক্ত অবশিষ্ট নেই।
\\
\\
\\
\\
\end{minipage}
\normalsize
\begin{minipage}[t]{0.6\linewidth}
\hspace{1em}
\end{minipage}
\begin{minipage}[t]{0.2\linewidth}
এআরও
\end{minipage}
\begin{minipage}[t]{0.2\linewidth}
আরও
\end{minipage}
\thispagestyle{laststyle}

\end{document}
